\renewcommand{\theposttitle}{Teaching, Breakdowns, True Wealth, \& Insights}
\renewcommand{\thepostdate}{2023-11-07}
\vspace*{\fill}
\begin{center}
\begin{minipage}{0.92\textwidth}
\small
Once you understand the relationship you’re presenting in a metaphor, you open the door for your imagination to create substitutions with the same concept.

For example, the iceberg metaphor is considered the most commonly used visual metaphor with the following relationship:
\end{minipage}
\end{center}
\vspace{4mm}
\begin{center}
\includegraphics[width=\textwidth, height=0.24\textheight, keepaspectratio]{C:/Users/James/Documents/GitHub/VisualIdeasBook/data/images/61c404c5e46dbc29a8bfe65ab06f98ca.png}
\end{center}
\vspace{2mm}
\begin{center}
\begin{minipage}{0.92\textwidth}
\small
Standard use-case for an iceberg metaphor

Once we understand the above surface vs below surface relationship, we can present it in a different context like this:
\end{minipage}
\end{center}
\vspace{5mm}
\begin{center}
\includegraphics[width=\textwidth, height=0.24\textheight, keepaspectratio]{C:/Users/James/Documents/GitHub/VisualIdeasBook/data/images/7fbb3f280fc0b100379774bb53d82efd.png}
\end{center}
\vspace{2mm}
\begin{center}
\begin{minipage}{0.92\textwidth}
\small
Can you think of a metaphor using this concept? \\
Shoot me a reply if you do. I’d love to hear it! 

And if you wanna learn more about how to create visual metaphors, \textbf{check out my course, Thinking in Visual Metaphors. }
\end{minipage}
\end{center}
\vspace*{\fill}